\chapter{Diagnostics and Failure Taxonomy}
\label{ch:bf-diagnostics}

Diagnostics are evidence about failures and geometry.  They are not proofs of
convergence.  BayesFilter should report diagnostics at three levels: filter
evaluation, derivative evaluation, and sampler behavior.

\section{Filter and Derivative Diagnostics}
\label{sec:bf-filter-derivative-diagnostics}

Filter diagnostics include factor pivots, reconstruction errors, solve
residuals, condition numbers, invalid-region counts, and finite-value checks.
Derivative diagnostics include score/Hessian finite checks, symmetry errors,
backend parity deltas, finite-difference residuals, and custom-gradient parity.
For SVD or eigendecomposition paths, singular-value or eigenvalue gap telemetry
is mandatory before HMC claims are made.

\section{Sampler Diagnostics}
\label{sec:bf-sampler-diagnostics}

Sampler diagnostics include acceptance rate, divergences, energy behavior,
effective sample size, split R-hat, mass-matrix conditioning, and tree depth if
NUTS is used diagnostically.  Short clean runs are smoke tests.  They can justify
the next experiment, but they do not establish convergence for an industrial
target.

\section{Warm-up and posterior precision}

A tuning allocation answers how much computation to spend learning geometry
and qualifying a frozen transition. It does not answer how long the resulting
chain must run to forget its initialization. BayesFilter therefore keeps all
verified $(L,\epsilon)$ pairs and asks the caller to choose a member before
posterior assessment. No R-hat, ESS or MCSE threshold changes tuning membership.

The selected member runs a separate discarded warm-up, followed by cumulative
retained draws. The shared keyword configuration inherits a minimum of 2,000
warm-up transitions per chain, a recent window of 1,000, modern R-hat at most
$1.05$, and a cap of 10,000. These numbers are operational owner-policy choices;
they are not estimates of a sufficient burn-in for every target. Stan's
windowed adaptation likewise allocates configured warm-up to fast step-size
and slow metric updates. Finishing that schedule is not a proof of stationarity.

\path{HMCPosteriorAssessmentPolicy} makes the additional requirements explicit:
minimum bulk and tail ESS during warm-up and retention, the number of
consecutive successful warm-up checks, and accuracy targets for the retained
estimators. Every eligible look reports rank-normalized split and folded R-hat,
bulk/tail ESS, original-scale mean ESS and mean MCSE. Custom callbacks add
requirements; they cannot replace core health and R-hat checks. The archived
wrapper retains its historical strict inequality and declared coordinate/ESS
screens, while using the same diagnostic arithmetic. The generic entry point
for a verified member is \path{run_hmc_posterior}.

Repeated successful windows are useful evidence of stability in the region
visited. Overlapping windows are not independent replications, and all chains
can remain in one unrepresentative mode. A small R-hat cannot detect a mode no
chain visits. Model-specific functionals, event indicators and log density can
be monitored with a named \path{quantities_fn} and a stable
\path{quantities_id}. If the warm-up cap is exhausted, the result reports
inconclusive equilibration and does not start retained inference. Warm-up draws
remain archived and are always excluded from posterior estimates.

For a scalar functional $g(\theta)$, let $X_{it}=g(\theta_{it})$ be draw $t$
of independent chain $i$, with $m$ chains of length $n$. Under stationarity and
an applicable central limit theorem, define the chain's long-run variance by
\begin{equation}
 \sigma_i^2=\gamma_i(0)+2\sum_{k\ge1}\gamma_i(k),\qquad
 \gamma_i(k)=\operatorname{Cov}(X_{i,t},X_{i,t+k}).
\end{equation}
Since the pooled estimate is $m^{-1}\sum_i\bar X_i$ and chains are independent,
its approximate Monte Carlo variance is
\begin{equation}
 \widehat{\operatorname{Var}}_{\rm MC}(\bar X)
   =\frac{1}{m^2 n}\sum_{i=1}^m\widehat\sigma_i^2,
 \qquad \operatorname{MCSE}(\bar X)
   =\sqrt{\widehat{\operatorname{Var}}_{\rm MC}(\bar X)}.
 \label{eq:bf-pooled-mean-mcse}
\end{equation}
Serial dependence enters through $\sigma_i^2$, while independence between
chains permits addition of their variances. Concatenating chains into a single
serial process would introduce artificial transitions between their endpoints.

For batch size $b$, let $a=\lfloor n/b\rfloor$ and let $\bar X_{ik}^{(b)}$
denote the mean in complete batch $k$. Ordinary batch means estimates
\begin{equation}
 \widehat\sigma_{i,b}^2=
 \frac{b}{a-1}\sum_{k=1}^{a}
   \left(\bar X_{ik}^{(b)}-a^{-1}\sum_{j=1}^{a}\bar X_{ij}^{(b)}\right)^2.
\end{equation}
The lugsail estimate combines two batch sizes
\citep[section 4, equation (7)]{vats2021lugsail}:
\begin{equation}
 \widehat\sigma_{i,L}^2=
 \frac{\widehat\sigma_{i,b}^2-c\widehat\sigma_{i,\lfloor b/r\rfloor}^2}{1-c},
 \qquad r\ge1,\quad 0\le c<1.
 \label{eq:bf-lugsail-mcse}
\end{equation}
For positively correlated chains, this construction can offset downward bias
in the ordinary estimate. It can also increase variability and yield a
nonpositive estimate. BayesFilter reports such estimates as unavailable
precision evidence instead of clipping them into a passing result. The
literature values $r=3,c=1/2$ and $b=\lfloor\sqrt n\rfloor$ are configurable
baselines, not universal recommendations. The minimum of twenty complete
batches is an operational stability floor, not a theorem about coverage.
Terminal incomplete batches are excluded from long-run variance estimation;
their count is recorded, and all retained draws still enter the posterior mean.

The autocorrelation option retains the identified TFP 0.25 positive-pairs ESS
calculation. Rank-normalized bulk ESS is useful for exploration diagnostics,
but original-scale mean uncertainty needs the ESS of the actual functional.
For a quantile, BayesFilter follows the indicator-ESS and order-statistic
construction in \citet[section 4.4]{vehtari2021rank}. A constant rare-event
indicator or a tied zero-width quantile interval supplies no reliable accuracy
evidence. Finite diagnostics do not establish the moments or mixing assumptions
needed by any of these estimators.

Each \path{HMCPrecisionTarget} specifies a mean or a quantile and an absolute
MCSE tolerance, an MCSE-to-posterior-SD tolerance, or both. Accuracy is assessed
on cumulative retained draws. No supplied tolerance means
\path{precision_not_requested}; a successful diagnostic screen then carries
no claim of sufficient estimator accuracy. At the retained cap, unmet targets
remain insufficient. Repeated MCSE screening is an operational stopping rule:
the implementation does not claim finite-sample or anytime-valid confidence
coverage. The asymptotic fixed-width results in
\citet[sections 2--3.1 and Appendix A]{flegalgongfixedwidth} require additional
regularity and estimator-consistency conditions. Lugsail estimates uncertainty
in retained means; it does not estimate how many burn-in draws are enough.

The diagnostic version \path{bayesfilter.hmc_diagnostic_math.v2} corrects two
prior arithmetic errors: normal-score ranks use $S+1/4$ in the denominator,
and modern R-hat is the square root of the variance ratio. Earlier reported
threshold results require recomputation from draws before current use.
Adjacent-epoch drift statistics remain descriptive: the old normalization
omits covariance between adjacent epoch means and is not a calibrated test.
The legacy Phase 29 warmup screen still compares these statistics with its
configured heuristic thresholds. Those comparisons can reject that legacy
screen; they are not part of the common posterior assessment and do not have
calibrated significance levels.

\lstinputlisting[language=Python,caption={A verified Gaussian member with explicit posterior accuracy targets.}]{examples/hmc_posterior_precision.py}

\section{Failure Taxonomy}
\label{sec:bf-failure-taxonomy}

Common failure classes are:
\begin{itemize}
  \item target-domain failures;
  \item factorization or solve failures;
  \item non-finite value or gradient failures;
  \item spectral derivative instability;
  \item JIT retracing or partial compilation;
  \item sampler geometry failures;
  \item unsupported missing-data or mask patterns.
\end{itemize}
Each failure should point to a next action: fix the model contract, change the
backend, add jitter or regularization, strengthen derivative validation, or
downgrade the backend to diagnostic use.

\section{Curvature Diagnostic Roles}
\label{sec:bf-curvature-diagnostic-roles}

Fixed-center curvature diagnostics have explicit roles:
\begin{description}
  \item[required gate] finite raw matrices, caller-declared raw inertia,
    selection-holdout error, projection burden, complete pairwise stability
    caps, and the independent audit error;
  \item[promotion veto] a failed required gate blocks nomination for an exact
    HMC mechanics canary;
  \item[repair trigger] failed one-factor holdout can trigger an identified
    two-factor fit, while unstable dense curvature triggers structured or
    consensus analysis;
  \item[explanatory] center score magnitude, optimizer iterations, descriptive
    runtimes, and any stability metric for which no cap was supplied.
\end{description}

SPD projection does not convert failed raw curvature into a pass.  Raw inertia
and matrices remain visible.  Likewise, a low score is not a curvature gate and
a stable nonzero-score center is not a MAP.  A locally useful geometry for an
exact Metropolis-corrected HMC canary remains unproved as a posterior covariance
or global-mode Hessian.

\texttt{geometry\_readiness\_blocked} rejects the current evidence and mass
handoff.  It does not imply that the model, target, or curvature-preconditioning
research direction has failed.
